\documentclass[11pt]{amsart}

\usepackage[margin=1.05in]{geometry}
\usepackage{amsmath,amssymb,amsthm,mathtools}
\usepackage{booktabs}
\usepackage{enumitem}
\usepackage{fancyvrb}
\usepackage[hidelinks]{hyperref}
\usepackage[capitalize,nameinlink]{cleveref}

\numberwithin{equation}{section}

\theoremstyle{plain}
\newtheorem{theorem}{Theorem}[section]
\newtheorem{proposition}[theorem]{Proposition}
\newtheorem{lemma}[theorem]{Lemma}
\newtheorem{corollary}[theorem]{Corollary}
\newtheorem{fact}{Fact}
\renewcommand{\thefact}{F\arabic{fact}}
\theoremstyle{definition}
\newtheorem{definition}[theorem]{Definition}
\newtheorem{remark}[theorem]{Remark}

\crefname{fact}{Fact}{Facts}

\newcommand{\Z}{\mathbb{Z}}
\newcommand{\Q}{\mathbb{Q}}
\newcommand{\F}{\mathbb{F}}
\newcommand{\Zq}{\Z[[q]]}
\newcommand{\coef}[1]{[q^{#1}]}
\newcommand{\one}{\mathbf 1}
\newcommand{\Jk}{\mathcal J}
\newcommand{\Sw}{\mathcal S}
\newcommand{\Tset}{\mathcal T}
\newcommand{\lab}{\operatorname{lab}}
\newcommand{\ol}{\overline}
\newcommand{\Leg}[2]{\Bigl(\frac{#1}{#2}\Bigr)}

\begin{document}

\title{Congruences modulo four for non-Rascoe partitions}

\author{Sumit Kumar Jha}
\address{Independent researcher, Greater Noida, Uttar Pradesh, India}
\email{sumitkumarjha75@gmail.com}

\subjclass[2020]{05A17, 11P83, 11F33}
\keywords{Non-Rascoe partitions; partition congruences; Rogers--Ramanujan
identities; Hecke--Rogers series; modular forms; involutions;
computer-assisted proofs}

\begin{abstract}
Let $b(n)$ be the number of partitions of $n$ into distinct parts in which the
number of parts does not occur as a part, and let $b_1(n)$ be the number of
those in which the number of parts plus one does not occur as a part. We prove
the conjecture of Dixit, Kumar and Srivastava that $b(29k+21)$ is divisible by
$4$ for every positive integer $k$ not divisible by $29$. More generally, if $p$
is a prime with $p\equiv21$ or $29\pmod{40}$, then $b(n)\equiv0\pmod4$ whenever
the exponent of $p$ in $40n+1$ is odd, and $b_1(n)\equiv0\pmod4$ whenever the
exponent of $p$ in $40n+9$ is odd. We also prove
$b(841n+21)+b(n)\equiv0\pmod4$ and $b_1(841n+189)+b_1(n)\equiv0\pmod4$ for
all $n\ge0$; the first of these was conjectured by Cuthbertson for $n\ge1$.

The proof reduces both generating functions modulo $4$ to a Hecke--Rogers
series and a count of lattice points in a cone. When the relevant
discriminant is not a square, the parity of the cone count is controlled by an
involution of Jackson, a sign built from the Legendre symbol, and a finite
cancellation identity that we check exhaustively by computer. When the
discriminant is a square, a divisor transform turns the problem into a
statement about a modular form of weight two, which is settled by Sturm's
theorem. The certificate data and the verification programs are included.
\end{abstract}

\maketitle
\setcounter{tocdepth}{1}
\tableofcontents

\section{Introduction}\label{sec:intro}

A partition of $n$ into distinct parts is a \emph{non-Rascoe partition} if its
number of parts is not itself one of the parts. Let $b(n)$ be the number of
non-Rascoe partitions of $n$, and let $b_1(n)$ be the number of partitions of
$n$ into distinct parts in which the number of parts plus one is not a part.
The empty partition gives $b(0)=b_1(0)=1$. For example, the partitions of $5$
into distinct parts are $5$, $4+1$ and $3+2$; the first contains its number of
parts, so $b(5)=2$. These functions were introduced by Dixit, Kumar and
Srivastava \cite{DKS}, who related them to Rogers--Ramanujan series and
conjectured that $b(29k+21)\equiv0\pmod4$ for $k\ge1$, $29\nmid k$
\cite[Conjecture~1]{DKS}. We prove this conjecture as a special case of the
following criterion. For a prime $p$ and a positive integer $D$, $v_p(D)$
denotes the exponent of $p$ in $D$.

\begin{theorem}\label{thm:main-valuation}
Let $p$ be a prime with $p\equiv21$ or $29\pmod{40}$, and let $n\ge0$. Then
\[
  v_p(40n+1)\ \text{odd}\ \Longrightarrow\ b(n)\equiv0\pmod4,\qquad
  v_p(40n+9)\ \text{odd}\ \Longrightarrow\ b_1(n)\equiv0\pmod4.
\]
In particular,
\begin{align*}
  b(29k+21)&\equiv0\pmod4&&(k\ge1,\ 29\nmid k),\\
  b_1(29k+15)&\equiv0\pmod4&&(k\ge0,\ k\not\equiv6\pmod{29}).
\end{align*}
\end{theorem}

For the prime $29$ we also obtain exact scaling relations, with no condition
on $n$.

\begin{theorem}\label{thm:main-scaling}
For every integer $n\ge0$,
\begin{align}
  b(841n+21)+b(n)&\equiv0\pmod4,\label{eq:scaling-b}\\
  b_1(841n+189)+b_1(n)&\equiv0\pmod4.\label{eq:scaling-b1}
\end{align}
\end{theorem}

For $n\ge1$, congruence \eqref{eq:scaling-b} is Cuthbertson's Conjecture~5.2
\cite{Cuthbertson}; the case $n=0$ also holds. Congruence \eqref{eq:scaling-b1}
is the companion statement for $b_1$. For every prime in \cref{thm:main-valuation},
\cref{cor:scaling-nonsquare} gives a similar scaling relation, provided the
associated number $40n+1$ or $40n+9$ is not a perfect square.

\subsection*{The cone count}
Both partition functions are controlled modulo $4$ by the following lattice
point count. For a positive integer $D\equiv1$ or $9\pmod{40}$ let
\[
  C(D)=\#\{(x,u,v)\in\Z_{>0}^3:\ 10u<x,\ x^2+40uv=D\},
\]
and put $C(D)=0$ if $D$ is not a positive integer. The equation forces
$x\equiv\pm1\pmod{10}$ if $D\equiv1\pmod{40}$ and $x\equiv\pm3\pmod{10}$ if
$D\equiv9\pmod{40}$.

\begin{theorem}\label{thm:cone-parity}
Let $p\equiv21$ or $29\pmod{40}$ be prime. If $D$ is a positive integer that is
not a perfect square, $D\equiv1$ or $9\pmod{40}$, and $p\mid D$, then
\begin{equation}\label{eq:cone-parity}
  C(D)\equiv C(D/p^2)\pmod2 .
\end{equation}
Consequently $C(D)$ is even whenever $v_p(D)$ is odd.
\end{theorem}

\subsection*{How the proof goes}
\Cref{sec:reduction} reduces the generating functions of $b$ and $b_1$ modulo $4$
to a Hecke--Rogers series, whose coefficients vanish at the relevant indices,
plus twice the generating function of $C$. So everything comes down to the
parity of $C(D)$.

For nonsquare $D$ (\cref{sec:involution,sec:certificate,sec:nonsquare}) we
enlarge the cone to the finite set of positive solutions of $x^2+2yz=D$. This
set carries Jackson's involution, a coordinate swap, and two rescalings, and a
sign built from the Legendre symbol modulo $p$ changes in a controlled way
under each of them. A finite identity, valid at every point of the set,
expresses the indicator of the cone as a sum of terms that cancel in pairs
under these maps. The identity depends only on residues modulo $20$ and $200$
and on finitely many linear inequalities, so it can be checked completely by
computer; the data and the checking program are part of the proof.

For square $D$, Jackson's involution has fixed boundaries and the argument
breaks down. Instead (\cref{sec:divisor,sec:recurrence,sec:scaling}) we write
$C(r^2)$ through a divisor sum, express that sum with Lambert series, and reduce
the problem at $p=29$ to the vanishing modulo $2$ of an explicit modular form of
weight two. Sturm's theorem reduces this to finitely many coefficients, which
are computed exactly.

\Cref{sec:verification} explains how to run the two computer verifications,
whose programs are printed in the appendices.

\section{Notation and standard facts}\label{sec:notation}

\subsection{Power series and arithmetic functions}
All power series have integer coefficients unless stated otherwise; $\Zq$ is the
ring of such series, and congruences between series are meant coefficient by
coefficient. For a series $F$, $\coef nF$ denotes the coefficient of $q^n$. We
write
\[
  (a;q)_r=\prod_{j=0}^{r-1}(1-aq^j)\quad(r\ge0),\qquad
  (a;q)_\infty=\prod_{j\ge0}(1-aq^j),
\]
with $(a;q)_0=1$. The \emph{odd part} of a series $F(q)$ is
$F_o(q)=\tfrac12\bigl(F(q)-F(-q)\bigr)$.

For an integer $N$, $\one_E$ is $1$ if the condition $E$ holds and $0$
otherwise. For positive integers $n$:
\begin{itemize}[leftmargin=2em]
\item $\tau(n)$ is the number of divisors of $n$, $\sigma(n)$ their sum, and
  $\mu$ is the M\"obius function;
\item $d_5(n)=\#\{d\mid n:5\nmid d\}$ and
  $\nu_5(n)=\#\{d\mid n:d\equiv2,3\pmod5\}$;
\item $(f*g)(n)=\sum_{d\mid n}f(d)g(n/d)$ is Dirichlet convolution.
\end{itemize}
An arithmetic function evaluated at a positive non-integer is $0$. For an odd
prime $p$, $\Leg ap$ is the Legendre symbol. Throughout, $p$ denotes a prime with
$p\equiv21$ or $29\pmod{40}$. For such $p$,
\begin{equation}\label{eq:legendre-values}
  \Leg{-1}p=1,\qquad \Leg2p=-1,\qquad \Leg5p=1,\qquad\text{hence}\qquad
  \Leg{10}p=\Leg{-2}p=\Leg{-10}p=-1 ,
\end{equation}
because $p\equiv5\pmod8$ and $p\equiv\pm1\pmod5$.

\subsection{Modular forms}
Write $q=e^{2\pi iz}$ with $z$ in the upper half-plane. For $L\ge1$,
\[
  \Gamma_0(L)=\Bigl\{\begin{pmatrix}a&b\\c&d\end{pmatrix}\in SL_2(\Z):L\mid c\Bigr\}.
\]
For a congruence subgroup $\Gamma$, $M_k(\Gamma)$ is the space of holomorphic
modular forms of integral weight $k$ on $\Gamma$ (holomorphic at all cusps). A
form of weight $k$ and character $\rho$ on $\Gamma_0(L)$ satisfies
$f\bigl(\frac{az+b}{cz+d}\bigr)=\rho(d)(cz+d)^kf(z)$. Dirichlet characters are
extended by $0$ to integers not coprime to their modulus, and $\ol\rho$ is the
complex conjugate character.

\subsection{Standard facts}

\begin{fact}[Dixit--Kumar--Srivastava {\cite[Thms.~1.2 and 4.6]{DKS}}]\label{fact:DKS}
For $\ell\in\{0,1\}$, with $b_0=b$,
\[
  \sum_{n\ge0}b_\ell(n)q^n=(-q;q)_\infty\sum_{r\ge0}\frac{(-1)^rq^{r(r+\ell)}}{(-q;q)_r}.
\]
\end{fact}

\begin{fact}[Rogers--Ramanujan--Bailey identities {\cite[Eqs.~(1.3)--(1.4)]{Sills}}]\label{fact:Sills}
As identities in $\Z[a][[q]]$,
\begin{align*}
  (aq;q)_\infty\sum_{r\ge0}\frac{a^rq^{r^2}}{(q;q)_r}
  &=1+\sum_{j\ge1}(-1)^ja^{2j}q^{j(5j-1)/2}\frac{(1-aq^{2j})(aq;q)_{j-1}}{(q;q)_j},\\
  (aq;q)_\infty\sum_{r\ge0}\frac{a^rq^{r(r+1)}}{(q;q)_r}
  &=\sum_{j\ge0}(-1)^ja^{2j}q^{j(5j+3)/2}\frac{(1-aq^{2j+1})(aq;q)_j}{(q;q)_j}.
\end{align*}
At $a=1$ these are the Rogers--Ramanujan identities.
\end{fact}

\begin{fact}[Hecke--Rogers identities, Andrews {\cite{Andrews}}; see also {\cite[\S1]{CGS}}]\label{fact:HR}
For $\ell\in\{0,1\}$,
\begin{equation}\label{eq:hecke-rogers}
  (q;q)_\infty\sum_{r\ge0}\frac{q^{r(r+\ell)}}{(-q;q)_r}
  =\sum_{\substack{x\in\Z,\ x\equiv2\ell+1\ (10)\\ j\in\Z,\ 10|j|<|x|}}
  \operatorname{sgn}(x)(-1)^jq^{(x^2-40j^2-(2\ell+1)^2)/40}.
\end{equation}
\end{fact}

\begin{fact}[Liouville-type identities {\cite[Thms.~5.6, 5.8, 5.10]{AAMW}}]\label{fact:Liouville}
Specialized to an even function as in \cref{prop:K5-reduction}, these give the
three evaluations \eqref{eq:liouville}.
\end{fact}

\begin{fact}[Weight-one Eisenstein series {\cite[\S5.3]{Stein}}]\label{fact:eisenstein}
For an odd primitive Dirichlet character $\rho$ of conductor $c$, the series
$E^\rho=\frac{L(0,\rho)}2+\sum_{n\ge1}\bigl(\sum_{d\mid n}\rho(d)\bigr)q^n$ lies in
$M_1(\Gamma_0(c),\rho)$. Here $L(s,\rho)=\sum_{n\ge1}\rho(n)n^{-s}$, continued
analytically to $s=0$.
\end{fact}

\begin{fact}[Theta series {\cite[\S\S2--3]{Buzzard}}]\label{fact:theta}
$\theta(z)=\sum_{a\in\Z}q^{a^2}$ is a modular form of weight $1/2$ on
$\Gamma_0(4)$, with the standard theta multiplier; the series $\theta(dz)$ are
modular on $\Gamma_0(4d)$.
\end{fact}

\begin{fact}[Sturm's theorem {\cite[Thm.~9.18]{Stein}}]\label{fact:sturm}
Let $\Gamma$ be a congruence subgroup of index $m$ in $SL_2(\Z)$ and let
$f\in M_k(\Gamma)$ have integer coefficients. If $\coef nf\equiv0\pmod2$ for all
$n\le km/12$, then $f\equiv0\pmod2$.
\end{fact}

\section{From partitions to the cone count}\label{sec:reduction}

For $\ell\in\{0,1\}$ put $b_0=b$ and define
\begin{align*}
  B_\ell(q)&=\sum_{n\ge0}b_\ell(n)q^n, &
  H_\ell(q)&=(q;q)_\infty\sum_{r\ge0}\frac{q^{r(r+\ell)}}{(-q;q)_r}=\sum_{n\ge0}h_\ell(n)q^n,\\
  \Theta_\ell(q)&=\sum_{m\in\Z}(-1)^mq^{m(5m+2\ell+1)/2}, &
  N_\ell(q)&=(q;q)_\infty\sum_{r\ge0}\frac{r\,q^{r(r+\ell)}}{(q;q)_r},
\end{align*}
and the Lambert series $\Lambda(q)=\sum_{j\ge1}q^j/(1-q^j)$.

Since $-1\equiv1\pmod2$, $(-q;q)_r\equiv(q;q)_r\pmod2$; so by the Rogers--Ramanujan
identities and the Jacobi triple product, $H_\ell\equiv\Theta_\ell\pmod2$. Also
\[
  \frac{(-q;q)_\infty}{(q;q)_\infty}=\prod_{j\ge1}\Bigl(1+\frac{2q^j}{1-q^j}\Bigr)
  \equiv1+2\Lambda\pmod4 .
\]
In the sum of \cref{fact:DKS} use $(-1)^r\equiv1+2r\pmod4$. The part with
$1$ gives $(-q;q)_\infty/(q;q)_\infty\cdot H_\ell\equiv(1+2\Lambda)H_\ell$, and in the
part with $2r$ we may replace $(-q;q)_\infty$ and $(-q;q)_r$ by $(q;q)_\infty$ and
$(q;q)_r$, since they agree modulo $2$. With $H_\ell\equiv\Theta_\ell\pmod2$ this gives
\begin{equation}\label{eq:B-mod4}
  B_\ell\equiv H_\ell+2\Theta_\ell\Lambda+2N_\ell\pmod4 .
\end{equation}

\subsection{The cone series}
Put
\[
  C_0(q)=\sum_{a\ge1}\Bigl(q^{a(5a-1)/2}\sum_{u=1}^{a-1}\frac{q^u}{1-q^u}
  +q^{a(5a+1)/2}\sum_{u=1}^{a}\frac{q^u}{1-q^u}\Bigr),
\]
and, with $\alpha_a=a(5a+3)/2$ and $e_a=(a+1)(5a+2)/2$,
\[
  C_1(q)=\sum_{a\ge0}\bigl(q^{\alpha_a}+q^{e_a}\bigr)\sum_{u=1}^{a}\frac{q^u}{1-q^u},\qquad
  U(q)=\sum_{a\ge0}q^{e_a}.
\]
An empty inner sum is $0$. Expanding the geometric series shows that
\begin{equation}\label{eq:cone-series}
  \coef nC_\ell(q)=C\bigl(40n+(2\ell+1)^2\bigr).
\end{equation}
Indeed, for $\ell=0$ the two sums correspond to $x=10a-1$ and $x=10a+1$, and for
$\ell=1$ to $x=10a+3$ and $x=10a+7$: one checks
$x^2-(2\ell+1)^2=40\cdot(\text{the leading exponent})$ in each case, the summation
variable $u$ ranges exactly over $10u<x$, and the geometric series in $u$
supplies $v$. The excluded value $x=1$ admits no $u$.

\begin{proposition}\label{prop:reduction}
One has
\[
  B_0\equiv H_0+2C_0\pmod4,\qquad B_1\equiv H_1+2C_1+2U\pmod4 .
\]
If $v_p(40n+(2\ell+1)^2)$ is odd, then
\begin{equation}\label{eq:b-equals-cone}
  b_\ell(n)\equiv2\,C\bigl(40n+(2\ell+1)^2\bigr)\pmod4 .
\end{equation}
\end{proposition}

\begin{proof}
Differentiate the identities of \cref{fact:Sills} with respect to $a$ and put
$a=1$; this is legitimate coefficientwise, since every coefficient is a
polynomial in $a$. Since $\frac{d}{da}(aq;q)_\infty\big|_{a=1}=-(q;q)_\infty\Lambda$
and, by the Rogers--Ramanujan identities and the Jacobi triple product,
$(q;q)_\infty\sum_rq^{r(r+\ell)}/(q;q)_r=\Theta_\ell$, the left sides become
$N_\ell-\Theta_\ell\Lambda$. On the right of the first identity, the derivative of
$a^{2j}$ is even and may be dropped modulo $2$. The derivative of $1-aq^{2j}$
gives $q^{2j}/(1-q^j)$, and the derivative of $(aq;q)_{j-1}$ gives
$(1-q^{2j})/(1-q^j)=1+q^j$ times $\sum_{u=1}^{j-1}q^u/(1-q^u)$, up to sign. Hence
\[
  N_0-\Theta_0\Lambda\equiv\sum_{j\ge1}q^{j(5j-1)/2}\Bigl((1+q^j)\sum_{u=1}^{j-1}\frac{q^u}{1-q^u}+\frac{q^{2j}}{1-q^j}\Bigr)
  \equiv C_0\pmod2 ;
\]
for the last step note $q^{j(5j-1)/2}\cdot q^j=q^{j(5j+1)/2}$ and
$q^{j(5j-1)/2}q^{2j}/(1-q^j)=q^{j(5j+1)/2}\,q^j/(1-q^j)$, which is the term $u=j$
of the second sum in $C_0$. In the second identity the derivative of
$1-aq^{2j+1}$ contributes $q^{e_j}$, and the finite product $(aq;q)_j$
contributes the summand of $C_1$; hence $N_1-\Theta_1\Lambda\equiv C_1+U\pmod2$.
Since $H_\ell\equiv\Theta_\ell\pmod 2$, substituting into \eqref{eq:B-mod4}
proves the two congruences.

For the last statement we use \eqref{eq:hecke-rogers}. If
$p\mid x^2-40j^2$ then, since $\Leg{10}p=-1$ by \eqref{eq:legendre-values},
$p$ divides both $x$ and $j$; dividing by $p^2$ repeatedly shows that every
nonzero integer of the form $x^2-40j^2$ has even $p$-adic valuation. So
$h_\ell(n)=0$ when $v_p(40n+(2\ell+1)^2)$ is odd. Finally
$40e_a+9=(10a+7)^2$, so $\coef nU=0$ when $v_p(40n+9)$ is odd. Together with
\eqref{eq:cone-series} this proves \eqref{eq:b-equals-cone}.
\end{proof}

\section{An involution and a sign}\label{sec:involution}

\subsection{The finite set and its maps}
Let $p$ be as in \cref{thm:cone-parity}, and let $D$ be a positive nonsquare with
$p\mid D$. First assume $D\equiv1\pmod{40}$, and put
\[
  \Tset_p(D)=\{(x,y,z)\in\Z_{>0}^3:\ x^2+2yz=D,\ \ p\nmid\gcd(x,y,z)\}.
\]
The last condition says only that $p$ does not divide all three coordinates.
The set is finite, its $x$-coordinates are odd, and $4\mid yz$.

Each triple $t=(x,y,z)$ falls into one of five \emph{branches}
$\beta(t)\in\{0,1,2,3,4\}$, according to the position of $2y$ relative to the
increasing numbers $x<2x+z<3x+2z<4x+4z$ (\cref{tab:jackson}). On branch $b$
\emph{Jackson's map} $\Jk$ is the linear map in the last column; see \cite{Jackson}
for the classical version, which is stated for prime discriminants.

\begin{table}[ht]
\caption{Branches of Jackson's map $\Jk$.}\label{tab:jackson}
\begin{tabular}{cll}
\toprule
$\beta$ & condition & $\Jk(x,y,z)$\\
\midrule
0 & $2y<x$ & $(x-2y,\ z+2x-2y,\ y)$\\
1 & $x<2y<2x+z$ & $(2y-x,\ y,\ 2x-2y+z)$\\
2 & $2x+z<2y<3x+2z$ & $(3x-2y+2z,\ 2x-y+2z,\ -2x+2y-z)$\\
3 & $3x+2z<2y<4x+4z$ & $(-3x+2y-2z,\ -2x+2y-z,\ 2x-y+2z)$\\
4 & $4x+4z<2y$ & $(x+2z,\ z,\ -2x+y-2z)$\\
\bottomrule
\end{tabular}
\end{table}

No triple lies on a boundary: $2y=x$ and $2y=3x+2z$ contradict the oddness of
$x$, while $2y=2x+z$ and $2y=4x+4z$ would give $D=(x+z)^2$ and $D=(x+2z)^2$. A
direct substitution shows that $\Jk$ preserves $x^2+2yz$, produces positive
coordinates, and sends branch $b$ to branch $(4,1,2,3,0)_b$. The linear maps on
branches $0$ and $4$ are inverse to each other, and those on branches $1,2,3$
are involutions; all have integral inverses, so $p\nmid\gcd$ is preserved. Hence
$\Jk$ is an involution of $\Tset_p(D)$.

The swap $\Sw(x,y,z)=(x,z,y)$ is a second involution. For $\ell\in\{2,5\}$
define the rational map
\[
  A_\ell(x,y,z)=(x,\ y/\ell,\ \ell z).
\]
It restricts to a bijection from the triples of $\Tset_p(D)$ with $\ell\mid y$ to
those with $\ell\mid z$; the condition $p\nmid\gcd$ is preserved because
$p\nmid10$.

\subsection{The sign}
Write $\chi(a)=\Leg ap$. By \eqref{eq:legendre-values},
$\chi(-1)=1$, $\chi(2)=-1$, $\chi(5)=1$. Let $\F_2=\Z/2\Z$. For
$t=(x,y,z)\in\Tset_p(D)$ define a sign $\epsilon(t)$ and a bit $g(t)\in\F_2$ by
\[
  \epsilon(x,y,z)=\begin{cases}\chi(y),&p\nmid y,\\ \chi(-2z),&p\mid y,\end{cases}
  \qquad (-1)^{g(t)}=\epsilon(t).
\]
This is well defined: if $p\mid y$ then $p\mid x$ (as $p\mid D$), so the
condition $p\nmid\gcd$ forces $p\nmid z$.

\begin{lemma}\label{lem:sign}
In $\F_2$, on the domains of the respective maps,
\[
  g(\Sw t)=g(t)+1,\qquad g(\Jk t)=g(t)+c_{\beta(t)},\quad (c_0,\dots,c_4)=(1,0,0,1,1),
\]
\[
  g(A_2t)=g(t)+1,\qquad g(A_5t)=g(t).
\]
\end{lemma}

\begin{proof}
Suppose first that $p\nmid y$. Since $p\mid D$, $z\equiv-x^2/(2y)\pmod p$. In
branches $0,2,3$ the new second coordinate is then congruent to
\[
  -\frac{(x-2y)^2}{2y},\qquad -\frac{(x-y)^2}{y},\qquad \frac{(x-2y)^2}{2y}
\]
respectively; in branch $1$ it is $y$, and in branch $4$ it is $z\equiv-x^2/(2y)$.
Since $\chi(-1)=1$ and $\chi(2)=-1$, these give the stated changes whenever the new
second coordinate is a unit modulo $p$. If it is divisible by $p$, use the new third
coordinate in the definition of $\epsilon$: in branches $0,2,3,4$ this gives
$\chi(-2y)$, $\chi(-y)$, $\chi(2y)$, $\chi(-2y)$ respectively (at the roots
$x\equiv2y$, $y$, $2y$, $0$), the same multipliers of $\chi(y)$.

If $p\mid y$, then $p\mid x$ and $p\nmid z$. The new second coordinates in
branches $0,2,3,4$ are congruent to $z,2z,-z,z$, and in branch $1$ the new third
coordinate is $z$; comparing with $\chi(-2z)$ gives the same table.

For $\Sw$: if $y,z$ are both units then $\chi(yz)=\chi(-x^2/2)=\chi(-2)=-1$, so
$\chi(z)=-\chi(y)$; if one of them is divisible by $p$, the second case of the
definition gives the same reversal. Finally $A_\ell$ multiplies the relevant
coordinate by $\chi(\ell)^{\pm1}$, giving $+1$ for $\ell=2$ and $0$ for $\ell=5$.
\end{proof}

\section{A pointwise cancellation identity}\label{sec:certificate}

The idea of this section is to write the indicator of the cone, viewed inside
$\Tset_p(D)$, as a sum of terms that cancel in pairs when summed over
$\Tset_p(D)$, because they are matched by $\Jk$, $\Sw$, $A_2$ and $A_5$. The
terms are given by a finite table (the \emph{certificate}); verifying the
identity is a finite computation, because everything involved depends only on
residues and on the position of the triple relative to finitely many lines.

\subsection{Labels}
For a positive rational triple $t=(x,y,z)$ put
\[
  L_0(t)=x-2y,\quad L_1(t)=2x-2y+z,\quad L_2(t)=3x-2y+2z,\quad L_3(t)=4x-2y+4z .
\]
These four numbers are strictly increasing. Extend $\beta$ to positive rational
triples: if none of the $L_i(t)$ is zero, $\beta(t)$ is the number of negative
$L_i(t)$; if $L_i(t)=0$, put $\beta(t)=5+i$. This agrees with the branches of
\cref{tab:jackson} away from boundaries. The \emph{label} of $t$ is
\[
  \lab(t)=\bigl(\beta(t),\ \beta(A_2t)\bigr).
\]
For $t\in\Tset_p(D)$ the first entry lies in $\{0,\dots,4\}$. The second entry
can also be $5$ or $7$, corresponding to
\begin{equation}\label{eq:boundary-labels}
  y=x\qquad\text{or}\qquad y=3x+4z ;
\end{equation}
the other two boundary values $6$ and $8$ would force $D=(x+2z)^2$ or
$D=(x+4z)^2$. So the labels that occur for integral triples lie in
\[
  \mathcal B=\{0,1,2,3,4\}\times\{0,1,2,3,4,5,7\}.
\]

\subsection{The certificate}
The certificate is printed in \cref{app:certificate}; the program there decodes
it to a data file
\[
  \texttt{original\_cone\_stratified.json}.
\]
It defines four functions
\[
  f_{\Jk},\ f_{\Sw},\ f_2,\ f_5:\ \mathcal B\times(\Z/20\Z)^3\times\F_2\longrightarrow\F_2 .
\]
Each of its $5485$ records $[\text{tag},[b,a],[r,s,u],g]$ marks one argument at
which the function takes the value $1$; all other values are $0$. The tags
\texttt{J}, \texttt{S}, \texttt{T}, \texttt{V} refer to $f_{\Jk}$, $f_{\Sw}$, $f_2$,
$f_5$. For a triple $t$ with bit $g(t)$ put
\[
  F_e(t)=f_e\bigl(\lab(t),\ t\bmod20,\ g(t)\bigr)\qquad(e\in\{\Jk,\Sw,2,5\}),
\]
where $t\bmod 20$ is taken componentwise.

\begin{lemma}\label{lem:cancellation}
For every $t=(x,y,z)\in\Tset_p(D)$ with $D\equiv1\pmod{40}$, in $\F_2$,
\begin{equation}\label{eq:cancellation}
\begin{aligned}
  \one_{\{2y<x,\ 5\mid y,\ 4\mid z\}}={}&F_{\Jk}(t)+F_{\Jk}(\Jk t)+F_{\Sw}(t)+F_{\Sw}(\Sw t)\\
  &+\one_{2\mid z}F_2(t)+\one_{2\mid y}F_2(A_2t)+\one_{5\mid z}F_5(t)+\one_{5\mid y}F_5(A_5t).
\end{aligned}
\end{equation}
\end{lemma}

\begin{proof}[Proof by exhaustive verification]
By \cref{lem:sign}, the bits of $\Jk t$, $\Sw t$, $A_2t$, $A_5t$ are determined by
$g(t)$ and $\beta(t)$. The remaining arithmetic data in \eqref{eq:cancellation}
depend only on
\[
  x\bmod20,\qquad y\bmod200,\qquad z\bmod20 :
\]
the modulus $200$ determines $y/2\bmod20$ when $2\mid y$ and $y/5\bmod20$ when
$5\mid y$, and these residues determine $x^2+2yz\bmod40$. The geometric data are
the five labels
\begin{equation}\label{eq:signature}
  \lab(t),\ \lab(\Jk t),\ \lab(\Sw t),\ \lab(A_2t),\ \lab(A_5t),
\end{equation}
which we call the \emph{signature} of $t$.

Normalize $x=1$. On each branch of $\Jk$, every inequality or equality defining
\eqref{eq:signature} is linear in $y,z$. The program
\texttt{verify\_original\_certificate.py} (\cref{app:verifier}) lists all
signatures that occur for some positive real $(y,z)$, by exact
Fourier--Motzkin elimination with strict and weak inequalities (an equality is
imposed as two weak inequalities). There are $263$ of them: $92$ without a
boundary label and $171$ with at least one. No approximation or sampling is used.

For an integral triple $s=(X,Y,Z)$ at which a term is evaluated, the program
imposes the necessary conditions $\lab(s)\in\mathcal B$,
$\beta(A_2s)=5\Rightarrow Y\equiv X\pmod{20}$ and
$\beta(A_2s)=7\Rightarrow Y\equiv3X+4Z\pmod{20}$, which follow from
\eqref{eq:boundary-labels}. It then runs over all residues $0\le x<20$,
$0\le y<200$, $0\le z<20$ with $x^2+2yz\equiv1\pmod{40}$, all signatures compatible
with these conditions, and both bits, and evaluates both sides of
\eqref{eq:cancellation} from the certificate. All $1\,156\,160$ identities hold,
$272\,960$ of them involving a boundary label. Every actual triple realizes one
of these combinations, which proves the lemma.
\end{proof}

\section{Proofs for nonsquare discriminants}\label{sec:nonsquare}

\begin{proof}[Proof of \cref{thm:cone-parity}]
First let $D\equiv1\pmod{40}$, and sum \eqref{eq:cancellation} over the finite
set $\Tset_p(D)$. The two $\Jk$-terms cancel because $\Jk$ permutes the set, and
the two $\Sw$-terms cancel for the same reason. For $\ell=2,5$, $A_\ell$ is a
bijection from $\{t:\ell\mid y\}$ onto $\{t:\ell\mid z\}$, so
$\sum_{\ell\mid z}F_\ell(t)=\sum_{\ell\mid y}F_\ell(A_\ell t)$ and these terms cancel too.
Hence the number of triples in $\Tset_p(D)$ with $2y<x$, $5\mid y$ and $4\mid z$ is
even.

If $D\equiv9\pmod{40}$, define $\Tset_p(D)$ by the same formula. For
$t\in\Tset_p(D)$ the triple $3t$ lies in $\Tset_p(9D)$ (as $p\ne3$), and
$9D\equiv1\pmod{40}$ is again a nonsquare divisible by $p$. Apply
\eqref{eq:cancellation} to $3t$ and sum over $t\in\Tset_p(D)$. All the maps
commute with multiplication by $3$, their domains are unchanged, and the
indicator on the left is unchanged because $3$ is prime to $20$. So the same
count is even in this case too.

Now put $y=5u$ and $z=4v$. The marked triples are exactly the triples counted by
$C(D)$ for which $p$ does not divide all of $x,u,v$. The triples with $p$
dividing all three are in bijection, by division by $p$, with those counted by
$C(D/p^2)$: the inequality $10u<x$ is preserved, and $p^2\equiv1\pmod{40}$ keeps
$D/p^2$ in the same class. (If $p^2\nmid D$ there are no such triples.) This
proves \eqref{eq:cone-parity}. Repeated division by $p^2$ reduces an odd
valuation to $1$, which gives the last assertion.
\end{proof}

\begin{proof}[Proof of \cref{thm:main-valuation}]
An odd valuation makes $40n+(2\ell+1)^2$ a nonsquare. Apply
\cref{prop:reduction} and \cref{thm:cone-parity}. For the two progressions use
\[
  40(29k+21)+1=29(40k+29),\qquad 40(29k+15)+9=29(40k+21);
\]
the valuation is exactly $1$ when $29\nmid k$, respectively when
$k\not\equiv6\pmod{29}$.
\end{proof}

\begin{corollary}\label{cor:scaling-nonsquare}
Let $p\equiv21$ or $29\pmod{40}$ be prime, let $\ell\in\{0,1\}$ and $n\ge0$, and
suppose $40n+(2\ell+1)^2$ is not a perfect square. Then
\[
  b_\ell\Bigl(p^2n+\frac{(2\ell+1)^2(p^2-1)}{40}\Bigr)\equiv s_p\,b_\ell(n)\pmod4,\qquad
  s_p=\begin{cases}1,&p\equiv1\pmod{10},\\-1,&p\equiv9\pmod{10}.\end{cases}
\]
\end{corollary}

\begin{proof}
The new index $n'$ satisfies $40n'+(2\ell+1)^2=p^2\bigl(40n+(2\ell+1)^2\bigr)$. In
\eqref{eq:hecke-rogers}, divisibility of $x^2-40j^2$ by $p$ forces $p\mid x$ and
$p\mid j$; substituting $x=s_ppX$, $j=pj'$ preserves the residue condition
$x\equiv2\ell+1\pmod{10}$ and the cone condition, multiplies the sign by $s_p$,
and leaves $(-1)^j$ unchanged since $p$ is odd. So $h_\ell(n')=s_ph_\ell(n)$.
\Cref{thm:cone-parity} gives $C$ at the new discriminant $\equiv C$ at the old one
modulo $2$. For $\ell=1$, both coefficients of $U$ vanish because both
discriminants are nonsquares. Now apply \cref{prop:reduction}.
\end{proof}

\begin{remark}
The nonsquare hypothesis is needed in general: numerically, the relation of
\cref{cor:scaling-nonsquare} fails at square discriminants for $p=61$ and
$p=109$. For $p=29$ it does hold at squares, which is the content of
\cref{thm:main-scaling} and \cref{sec:divisor,sec:recurrence,sec:scaling}.
\end{remark}

\section{Square discriminants: a divisor transform}\label{sec:divisor}

The nonsquare hypothesis in \cref{thm:cone-parity} was used to keep triples off
the boundaries of $\Jk$. For square discriminants we replace the cone
inequality by a divisor transform and then work with Lambert series.

\subsection{The divisor transform}
For a positive odd integer $m$ prime to $5$ put $\gamma(m)=C(m^2)$. For $d\ge1$ let
$\mathcal L_d(q)=q^d/(1-q^d)$, and define
\[
  \mathcal A(q)=\sum_{\substack{a,s\ge1\\5\mid a}}q^{2as}\bigl(\mathcal L_a(q)\mathcal L_{2s}(q)+\mathcal L_{2a}(q)\mathcal L_s(q)\bigr),
  \qquad A_m=\coef m\mathcal A(q).
\]
Only finitely many terms contribute to each coefficient.

\begin{proposition}\label{prop:divisor-transform}
For every positive odd integer $m$ prime to $5$,
\[
  A_m=\sum_{d\mid m}\gamma(d),\qquad \gamma(m)=\sum_{d\mid m}\mu(d)A_{m/d}.
\]
\end{proposition}

\begin{proof}
For a triple counted by $\gamma(m)$ put $a=(m-x)/2$ and $b=(m+x)/2$; the positive
integer matrix
\[
  T=\begin{pmatrix}a&5u\\2v&b\end{pmatrix}
\]
has determinant $ab-10uv=(m^2-x^2-40uv)/4=0$ and trace $m$. Let $g$ be the gcd of
its entries. Then $g\mid m$, so $g$ is prime to $10$ and divides $x,u,v$.
Conversely $d=\gcd(x,u,v)$ satisfies $d^2\mid m^2$, so $d\mid m$, and since $d$ is
odd it divides $a$ and $b$. Thus $g=\gcd(x,u,v)$.

A rank-one positive integer matrix factors as a column $(h,k)^{T}$ times a row
$(r,s)$ in exactly $\tau(g)$ ways: dividing the first row by its gcd gives the
unique factorization with $\gcd(r,s)=1$, whose column has gcd $g$, and every other
factorization moves a divisor of $g$ from the column to the row. These
factorizations correspond to positive quadruples $(h,k,r,s)$ with
\[
  hr+ks=m,\qquad 5\mid hs,\qquad 2\mid kr,\qquad ks-hr>2hs,
\]
with inverse $x=ks-hr$, $u=hs/5$, $v=kr/2$. Write $k=t+2h$. The inequality forces
$t>0$ and becomes $ts>hr$; the trace condition becomes $hr+ts+2hs=m$, and the parity
condition $2\mid tr$. Without the inequality, the involution
$(h,t,r,s)\mapsto(s,r,t,h)$ exchanges $ts$ and $hr$, and equality $ts=hr$ would make
$m$ even. So the count without the inequality is twice the sum of
$\tau(\gcd(x,u,v))$ over the triples counted by $\gamma(m)$.

For fixed $h,s$, inclusion--exclusion for $2\mid tr$ gives
$\mathcal L_h\mathcal L_{2s}+\mathcal L_{2h}\mathcal L_s-\mathcal L_{2h}\mathcal L_{2s}$; after
multiplication by $q^{2hs}$ the last product has only even exponents. The other two
terms are symmetric in $h,s$, so inclusion--exclusion for $5\mid hs$ makes their
coefficient at $q^m$ twice the coefficient with $5\mid h$ (the case $5\mid h$ and $5\mid s$
contributes only multiples of $5$). Hence
$A_m=\sum\tau(\gcd(x,u,v))$ over the triples counted by $\gamma(m)$. Expanding
$\tau$ and dividing each triple by a common divisor $d\mid m$ gives
$A_m=\sum_{d\mid m}\gamma(m/d)$, and M\"obius inversion gives the second formula.
\end{proof}

\subsection{Products of Lambert series}
Unspecified summation indices are positive integers. Put
\[
  \mathcal O_d(q)=\frac{q^d}{1-q^{2d}},\quad
  \mathcal O=\sum_{d\ \mathrm{odd}}\mathcal O_d,\quad
  \mathcal O_5=\sum_{\substack{d\ \mathrm{odd}\\5\mid d}}\mathcal O_d,\quad
  \mathcal V=\sum_{d}\mathcal L_{2d},\quad
  \mathcal V_5=\sum_{d}\mathcal L_{10d},
\]
and, for $1\le j\le10$ and $a\in\{1,\dots,10\}$,
\[
  \mathcal P_j(q)=\sum_{r\ \mathrm{odd}}\frac{q^{jr}}{1-q^{10r}}=\sum_{d\equiv j\ (10)}\mathcal O_d(q),
  \qquad
  \mathcal T_a(q)=\sum_{d\equiv a\ (10)}\mathcal L_{2d}(q),
\]
where $a=10$ stands for the class $0$. (The second expression for $\mathcal P_j$
follows by expanding and exchanging the two indices.) In particular
$\mathcal O_5=\mathcal P_5$. Let $\mathcal A_{\mathrm{odd}}=\sum_{N\ \mathrm{odd}}A_Nq^N$, and
\[
  \mathcal F_5=\sum_{a\in\{1,3,5,7,9\}}\mathcal P_a\mathcal T_{10-a},\qquad
  \mathcal G_5=\sum_{j=1}^5\mathcal P_j\mathcal P_{j+5}.
\]

\begin{proposition}\label{prop:lambert}
In $\Zq$,
$\mathcal A_{\mathrm{odd}}=\mathcal O_5\mathcal V+\mathcal O\mathcal V_5-\mathcal F_5-\mathcal G_5$.
\end{proposition}

\begin{proof}
Extracting odd powers in the definition of $\mathcal A$ gives
\[
  \mathcal A_{\mathrm{odd}}=\sum_{m\ \mathrm{odd},\ s\ge1}\bigl(\one_{5\mid m}+\one_{5\mid s}\bigr)q^{2ms}\mathcal O_m\mathcal L_{2s}
\]
(in the second product of $\mathcal A$, exchange the names of the two indices
after extracting odd powers; terms with both divisibilities have multiplicity
two). Call the two parts $D_5$ and $E_5$. Expanding $\mathcal O_m$ and exchanging
odd factors gives
\[
  \sum_{m\ \mathrm{odd}}q^{2ms}\mathcal O_m=\sum_{\substack{r\ \mathrm{odd}\\r>2s}}\mathcal O_r,\qquad
  \sum_{\substack{m\ \mathrm{odd}\\5\mid m}}q^{2ms}\mathcal O_m=\sum_{\substack{r\ \mathrm{odd}\\r>2s}}\mathcal O_{5r}.
\]
Removing the finite initial segments,
\[
  D_5=\mathcal O_5\mathcal V-\sum_{r,t\ \mathrm{odd}}\mathcal O_{5r}\mathcal L_{r+t},\qquad
  E_5=\mathcal O\mathcal V_5-\sum_{5\mid s}\mathcal L_{2s}\sum_{\substack{r\ \mathrm{odd}\\r<2s}}\mathcal O_r ,
\]
where in the first we parametrized the removed segment by $t=2s-r$.

For the second removed sum use the identity
$\mathcal O_u\mathcal L_{2v}+\mathcal O_v\mathcal L_{2u}=(\mathcal O_u+\mathcal O_v)\mathcal L_{u+v}$ and
pair $(u,v)$ with $(v,u)$ (a diagonal term is $\mathcal O_u\mathcal L_{2u}$):
\[
  \sum_{5\mid s}\mathcal L_{2s}\sum_{\substack{u\ \mathrm{odd}\\u<2s}}\mathcal O_u
  =\sum_{\substack{u,v\ \mathrm{odd}\\5\mid u+v}}\mathcal O_u\mathcal L_{2v}
  =\sum_{a\in\{1,3,5,7,9\}}\mathcal P_a\mathcal T_{10-a}=\mathcal F_5,
\]
since for odd $u,v$ the condition $5\mid u+v$ means $v\equiv10-a\pmod{10}$ when
$u\equiv a$. For the first removed sum the same pairing, with $a=q^r$, $b=q^t$ and
$(1-(ab)^5)/(1-ab)=1+ab+\dots+(ab)^4$, gives
\[
  (\mathcal O_{5r}+\mathcal O_{5t})\mathcal L_{r+t}=\sum_{j=1}^5\frac{q^{(j+5)r+jt}+q^{jr+(j+5)t}}{(1-q^{10r})(1-q^{10t})},
\]
and the diagonal terms also match:
$\sum_{j=1}^5q^{(2j+5)r}/(1-q^{10r})^2=q^{7r}/((1-q^{2r})(1-q^{10r}))=\mathcal O_{5r}\mathcal L_{2r}$.
Thus $\sum_{r,t\ \mathrm{odd}}\mathcal O_{5r}\mathcal L_{r+t}=\sum_{j=1}^5\mathcal P_j\mathcal P_{j+5}=\mathcal G_5$.
All rearrangements are coefficientwise finite.
\end{proof}

\subsection{Reduction to one divisor convolution}
For a positive odd integer $N$ put
\[
  K_5(N)=\sum_{t=1}^{(N-1)/2}d_5(N-2t)\,d_5(t),\qquad
  R(N)=\coef N(\mathcal O_5\mathcal V+\mathcal O\mathcal V_5),
\]
and, for $N$ also prime to $5$,
\[
  Q(N)=\sum_{d\mid N}q_0(d),\qquad
  q_0(d)=\frac{d-1}2+2\Bigl\lfloor\frac d5\Bigr\rfloor-3\Bigl\lfloor\frac d{10}\Bigr\rfloor .
\]

\begin{proposition}\label{prop:K5-reduction}
For every positive odd $N$ prime to $5$, $K_5(N)-Q(N)$ is even and
\begin{equation}\label{eq:A-mod2}
  A_N\equiv R(N)+\frac{K_5(N)-Q(N)}2\pmod2 .
\end{equation}
\end{proposition}

\begin{proof}
Extend $\mathcal T_a$ to all ten classes as above. Since $\mathcal O_d=\mathcal L_d-\mathcal L_{2d}$,
direct expansion gives $\mathcal P_{2j}=\mathcal T_j+\mathcal T_{j+5}-\mathcal T_{2j}$ for
$1\le j\le5$. Using this, $\mathcal F_5+\mathcal G_5$ is congruent modulo $2$ to the
generating function of the positive quadruples $(m,r,s,t)$ with $m,r$ odd, counted
with weight $w(m,s)$ at the exponent $mr+2st$, where $w$ is given by
\cref{tab:weights}. (The term $\mathcal P_5\mathcal T_5$ occurs twice and drops out.)

\begin{table}[ht]
\caption{The nonzero weights $w(m,s)$; all other weights are $0$.}\label{tab:weights}
\begin{tabular}{cc}
\toprule
$m\bmod10$ & $s\bmod10$ with $w(m,s)=1$\\
\midrule
1 & 3, 6, 8, 9\\
3 & 4, 7, 8, 9\\
7 & 1, 2, 3, 6\\
9 & 1, 2, 4, 7\\
\bottomrule
\end{tabular}
\end{table}

Write $\delta(z)=\one_{5\mid z}$ and $\Delta(a,b)=\delta(a-b)-\delta(a+b)$. Then, as an
identity of integers,
\[
  2w(m,s)=\one_{5\nmid ms}+\Delta(m,2s)+(-1)^s\Delta(m,s),
\]
which is checked on the residue classes modulo $10$ (when $5\mid ms$ both differences
vanish). Let $W(N)$ be the sum of $w(m,s)$ over $mr+2st=N$ with $m,r,s,t>0$ and $m,r$
odd. \Cref{prop:lambert} gives $A_N\equiv R(N)+W(N)\pmod2$.

Let $S_6,S_8,S_{10}$ be the sums of $\Delta(a,b)$ over positive $a,b,x,y$ with
$ax+by=N$, under the respective extra conditions: $a,x,y$ odd; $a,x$ odd; $a,b,x$
odd. Put
\[
  J(N)=\frac{\sigma(N)-\tau(N)}2,\qquad \Phi_j(N)=\sum_{d\mid N}\Bigl\lfloor\frac dj\Bigr\rfloor .
\]
Specializing the Liouville-type identities of \cref{fact:Liouville} to the even
function $\delta$ gives
\begin{equation}\label{eq:liouville}
  S_6=0,\qquad S_8=J(N)-\Phi_5(N),\qquad S_{10}=J(N)-\Phi_{10}(N);
\end{equation}
here $\delta(0)=1$ while $\delta(d)=0$ for every $d\mid N$, and the remaining inner
sums count multiples of $5$ or even multiples of $5$.

Let $X_2$ be the sum of $\Delta(m,2s)$ over $mr+2st=N$ with $m,r$ odd, and let
$X_{\mathrm{odd}}$, $X_{\mathrm{even}}$ be the sums of $\Delta(m,s)$ over the same set with
$s$ odd, respectively even. Splitting $S_8$ by the parity of $b$ and then of $y$
gives
\[
  X_2=S_8-S_{10},\qquad X_{\mathrm{odd}}=S_{10},\qquad X_{\mathrm{even}}=S_8-S_6-S_{10}.
\]
The sum of $\one_{5\nmid ms}$ is exactly $K_5(N)$. Summing the expression for $2w$
therefore gives the integer identity
\[
  2W(N)=K_5(N)+X_2+X_{\mathrm{even}}-X_{\mathrm{odd}}=K_5(N)+2S_8-S_6-3S_{10}=K_5(N)-Q(N),
\]
because $Q(N)=J(N)+2\Phi_5(N)-3\Phi_{10}(N)$. So $K_5(N)-Q(N)$ is even, and
substituting $W(N)=(K_5(N)-Q(N))/2$ proves \eqref{eq:A-mod2}.
\end{proof}

\section{A recurrence for the convolution at 29}\label{sec:recurrence}

Throughout this section $p=29$. For a positive integer $N$ divisible by $p$ write
$N=p^em$ with $e=v_p(N)\ge1$ and $p\nmid m$.

\begin{theorem}\label{thm:K5-recurrence}
For every positive odd integer $N$ divisible by $p$,
\[
  K_5(N)-K_5(N/p^2)\equiv2\bigl(\nu_5(m)+d_5(m)\bigr)\pmod4,\qquad m=N/p^{v_p(N)} .
\]
\end{theorem}

The proof builds an integral modular form whose coefficients at multiples of $p$
measure the failure of this congruence, and shows that it vanishes modulo $2$ by
Sturm's theorem.

\subsection{Characters, theta series and Eisenstein series}
Let $\psi$ be the quartic character modulo $5$ with
$\psi(1)=1$, $\psi(2)=i$, $\psi(3)=-i$, $\psi(4)=-1$. Let $\chi_4(n)=(-1)^{(n-1)/2}$ for
odd $n$ (and $0$ for even $n$), and let $\chi_8$ be the character modulo $8$ with
$\chi_8(1)=\chi_8(7)=1$, $\chi_8(3)=\chi_8(5)=-1$. Let $\lambda$ be the quartic character
modulo $p$ with $\lambda(2)=i$; since $2$ has order $28$ modulo $29$, this determines
$\lambda$, and
\[
  \lambda(-1)=-1,\qquad \lambda(5)=-1,\qquad \lambda^2(a)=\Leg a{29}.
\]
For $d\ge1$ define the theta series
\[
  \theta_d=\sum_{a\in\Z}q^{da^2},\qquad
  \theta^+_d=\sum_{a\ge1}q^{da^2},\qquad
  \theta^{\mathrm{odd}}_d=\sum_{a\ge1,\ a\ \mathrm{odd}}q^{da^2}=\frac{\theta_d-\theta_{4d}}2 .
\]
For $\rho\in\{\psi,\ol\psi,\lambda,\ol\lambda\}$ let $E^\rho$ be the weight-one
Eisenstein series of \cref{fact:eisenstein} (each $\rho$ is odd and primitive), and put
\[
  E_R=\frac{E^\lambda+E^{\ol\lambda}}2,\qquad E_I=\frac{E^\psi-E^{\ol\psi}}{2i},
\]
taking real and imaginary parts of the coefficients.

\subsection{A quartic-character lift}
Put
\[
  E(q)=E_R(q)+E_R(q^5)-E_R(q^p)-E_R(q^{5p}).
\]
Its constant term is $0$. Since $\lambda(p)=0$ and $\lambda(5)=-1$, its coefficients are
the integers
\[
  \coef nE=\begin{cases}\operatorname{Re}\sum_{d\mid n,\ 5\nmid d}\lambda(d),&p\nmid n,\\ 0,&p\mid n.\end{cases}
\]
If $\Leg n{29}=-1$, then $d_5(n)$ is even and
\begin{equation}\label{eq:d5-half}
  \frac{d_5(n)}2\equiv\coef nE\pmod2 .
\end{equation}
Indeed, in each pair of complementary divisors $d,n/d$ exactly one value of
$\lambda$ is real ($\pm1$) and the other imaginary, since
$\lambda^2(d)\lambda^2(n/d)=\Leg n{29}=-1$. So $\tau(n)$ is even and $\tau(n)/2$ has the
parity of $\operatorname{Re}\sum_{d\mid n}\lambda(d)$. If $5\mid n$, apply this also to
$n/5$ (still a nonresidue, as $\Leg5{29}=1$) and subtract, using
$d_5(n)=\tau(n)-\tau(n/5)$.

Write $F_o$ for the odd part of a series (\cref{sec:notation}). Define
\[
  \Omega_d=\theta^{\mathrm{odd}}_d+\theta^{\mathrm{odd}}_{5d},\qquad
  \vartheta_d=\frac{\theta_{2d}-\theta_{10d}}2=\theta^+_{2d}-\theta^+_{10d},\qquad
  X=\Omega_1-\Omega_{p^2},
\]
all with integer coefficients, and
\[
  \mathcal R=\Omega_{p^2}\vartheta_p+\Omega_p\vartheta_{p^2},\qquad
  \mathcal M=X\,E(q^2)+(\vartheta_1-\vartheta_{p^2})E_o+\mathcal R .
\]

\subsection{The convolution difference}
We claim that for every positive odd $N$ divisible by $p$,
\begin{equation}\label{eq:conv-diff}
  \frac{K_5(N)-K_5(N/p^2)}2\equiv\coef N\mathcal M\pmod2 ;
\end{equation}
in particular the numerator is even. To avoid dividing a congruence by $2$, let
\[
  \mathcal D=\sum_{n\ge1}d_5(n)q^n,\qquad
  \mathcal D_0=\sum_{p\nmid n}d_5(n)q^n,\qquad
  \mathcal Z=\sum_{p\mid n}d_5\bigl(n/p^{v_p(n)}\bigr)q^n .
\]
Since $d_5(p^em)=(e+1)d_5(m)$ for $p\nmid m$, we have the exact identity
$\mathcal D(q)-\mathcal D(q^{p^2})=\mathcal D_0(q)+2\mathcal Z(q)$.

The generating function of $K_5(N)$ at odd $N$ is $\mathcal D_o(q)\mathcal D(q^2)$.
Expand the difference between it and its image under $q\mapsto q^{p^2}$ using the
identity just stated. At exponents divisible by $p$, a product of a factor
supported on multiples of $p$ with one supported away from multiples of $p$
contributes nothing. So, modulo $4$, at such exponents the difference has the
same coefficients as
\[
  \mathcal D_{0,o}(q)\mathcal D_0(q^2)+2\mathcal Z_o(q)\mathcal D(q^{2p^2})+2\mathcal D_o(q^{p^2})\mathcal Z(q^2).
\]
In the first product, no term at an exponent divisible by $p$ has both divisor
counts odd: $d_5(n)$ is odd only when $n$ is a square or five times a square, so
such an exponent would be one of $x^2+2y^2$, $x^2+10y^2$, $5x^2+2y^2$,
$5x^2+10y^2$ with $p\nmid xy$, which is impossible because
$\Leg{-2}p=\Leg{-10}p=-1$. So its coefficient is even, and modulo $4$ only terms
with exactly one odd divisor count contribute. The odd factor is supported on
squares or five times squares; the argument of the other factor is then a
nonresidue modulo $p$, so \eqref{eq:d5-half} applies to it. This shows that half
of the first product, at the stated exponents, has the parity of
$\coef N\bigl(X\,E(q^2)+(\vartheta_1-\vartheta_{p^2})E_o\bigr)$; here
$\theta^+_{2d}+\theta^+_{10d}\equiv\vartheta_d\pmod2$ accounts for the theta differences.

For the other two products, pairing complementary divisors gives
\[
  \mathcal D\equiv\theta^+_1+\theta^+_5,\qquad
  \mathcal Z\equiv\theta^+_p+\theta^+_{5p}+\theta^+_{p^2}+\theta^+_{5p^2}\pmod2 ,
\]
the second separating odd and positive even valuations at $p$. Taking odd parts
and substituting, the two equal terms $\Omega_{p^2}(\theta^+_{2p^2}+\theta^+_{10p^2})$
cancel modulo $2$, and what remains is $\mathcal R$. This proves \eqref{eq:conv-diff}.

\subsection{An integral form of weight two}
Put
\[
  \mathcal N=(E_I)_o(q^p)+(E_I)_o(q^{p^2}),\qquad \mathcal Y=\Omega_p+\Omega_{p^2}.
\]
For odd $N=p^em$ with $p\nmid m$,
\[
  \coef N\mathcal N\equiv\nu_5(m),\qquad \coef N\mathcal Y\equiv d_5(m)\pmod2 .
\]
For the first: $\coef nE_I\equiv\nu_5(n)\pmod2$, and
$\nu_5(p^am)=(a+1)\nu_5(m)$ because multiplication by $p\equiv4\pmod5$ permutes the
two classes $2,3$ modulo $5$; the two shifted series contribute
$e\,\nu_5(m)+(e-1)\nu_5(m)$ (the second term absent for $e=1$). The second again
separates odd and positive even valuations at $p$. Define
\[
  \mathfrak F=\theta_{40}\bigl(X\,E(q^2)+(\vartheta_1-\vartheta_{p^2})E_o\bigr)
  +\theta_{40}^2(\mathcal R-\mathcal N)-\theta_{40}^3\,\mathcal Y .
\]
Every coefficient of $\mathfrak F$ is an integer: the constant term of $E$ is $0$,
taking the odd part removes the constant term of $E_I$, and the theta differences
$\vartheta_d$ have integer coefficients. Since $\theta_{40}\equiv1\pmod2$,
\begin{equation}\label{eq:F-coefficients}
  \coef N\mathfrak F\equiv\frac{K_5(N)-K_5(N/p^2)}2-\nu_5(m)-d_5(m)\pmod2
  \qquad(p\mid N,\ N\ \mathrm{odd}).
\end{equation}

\subsection{Modularity and the Sturm bound}
Define the even characters
\[
  \kappa_1=\chi_4\lambda,\qquad \kappa_2=\chi_4\psi,\qquad \kappa_3=\chi_8,
\]
of orders $4$, $4$ and $2$. Each summand of $\mathfrak F$ is a sum of weight-two
products of theta series (\cref{fact:theta}) and Eisenstein series
(\cref{fact:eisenstein}), holomorphic at every cusp, of level dividing $160p^2$ and
with character in $\langle\kappa_1,\kappa_2,\kappa_3\rangle$; \cref{tab:characters}
lists the characters ($a,b\in\{0,1\}$ independently).

\begin{table}[ht]
\caption{Characters of the weight-two constituents of $\mathfrak F$.}\label{tab:characters}
\begin{tabular}{ll}
\toprule
summand & possible characters\\
\midrule
$\theta_{40}X\,E(q^2)$ & $\kappa_1^{\pm1}\kappa_3\kappa_2^{2a}$\\
$\theta_{40}(\vartheta_1-\vartheta_{p^2})E_o$ & $\kappa_1^{\pm1}\kappa_2^{2a}$\\
$\theta_{40}^2\mathcal R$ & $\kappa_1^2\kappa_3\kappa_2^{2a}$\\
$\theta_{40}^2\mathcal N$ & $\kappa_2^{\pm1}$\\
$\theta_{40}^3\mathcal Y$ & $\kappa_1^{2a}\kappa_3\kappa_2^{2b}$\\
\bottomrule
\end{tabular}
\end{table}

Taking the odd part of a weight-one form of odd level $L$ gives a form of level
$4L$ with the same character (write the projection using translation by $1/2$
and conjugate on $\Gamma_0(4L)$), so the odd projections stay within the stated
level bound.

For a series $f=\sum a(n)q^n$ let $U_pf=\sum a(pn)q^n$; analytically
$U_pf(z)=\frac1p\sum_{j=0}^{p-1}f\bigl(\frac{z+j}p\bigr)$. If $f$ has weight two, level
$160p^2$ and one of the characters above, then $U_pf$ has level $160p$ and the same
character. To see this, let $\gamma=\begin{pmatrix}a&b\\c&d\end{pmatrix}\in\Gamma_0(160p)$,
$M_j=\begin{pmatrix}1&j\\0&p\end{pmatrix}$, and choose $j'$ modulo $p$ with
$j'\equiv a^{-1}(b+jd)\pmod p$. Then
\[
  M_j\gamma=\begin{pmatrix}a+jc&(b+jd-(a+jc)j')/p\\ pc&d-cj'\end{pmatrix}M_{j'} .
\]
The first matrix lies in $\Gamma_0(160p^2)$, its lower-right entry is congruent to
$d$ modulo $40p$, and $j\mapsto j'$ is a permutation. Summing the transformation
laws over $j$ proves the claim, because every character involved has conductor
dividing $40p$; holomorphy at the cusps is preserved by this finite sum.

Consequently $U_p\mathfrak F\in M_2(\Gamma_{29})\cap\Zq$, where
\[
  \Gamma_{29}=\Bigl\{\begin{pmatrix}a&b\\c&d\end{pmatrix}\in\Gamma_0(160p):
  \kappa_1(d)=\kappa_2(d)=\kappa_3(d)=1\Bigr\}.
\]
The character map $d\mapsto(\kappa_1(d),\kappa_2(d),\kappa_3(d))$ has image of order
$4\cdot4\cdot2=32$: by the Chinese remainder theorem the residues modulo $p$ and $5$
prescribe $\kappa_1,\kappa_2$ independently, and the residue modulo $8$ prescribes
$\kappa_3$. Hence
\[
  [SL_2(\Z):\Gamma_{29}]=32\cdot160p\Bigl(1+\frac12\Bigr)\Bigl(1+\frac15\Bigr)\Bigl(1+\frac1{29}\Bigr)=276480,
\]
and Sturm's bound (\cref{fact:sturm}) is $2\cdot276480/12=46080$.

\begin{proof}[Proof of \cref{thm:K5-recurrence}]
The program \texttt{prove\_square\_scaling\_29.py} (\cref{app:sturm}) computes
$\mathfrak F$ modulo $2$ through degree $29\cdot46080=1\,336\,320$, using Python
integers as bit vectors for polynomials over $\F_2$ (addition is exclusive-or,
multiplication is shift-and-xor), so no rounding occurs. It checks that every
coefficient of $U_{29}\mathfrak F$ through degree $46080$ is $0$ modulo $2$; the
constant term is $0$, and all even coefficients vanish because $\mathfrak F$ is
supported on odd exponents. By \cref{fact:sturm}, $U_{29}\mathfrak F\equiv0\pmod2$.
Now \eqref{eq:F-coefficients} gives the theorem.
\end{proof}

\section{The scaling congruences}\label{sec:scaling}

We keep $p=29$ and transfer \cref{thm:K5-recurrence} to the divisor transform of
\cref{prop:divisor-transform}.

\begin{proposition}\label{prop:A-recurrence}
For every positive odd $N$ prime to $5$ and divisible by $p$,
\[
  A_N-A_{N/p^2}\equiv\nu_5(m)\pmod2,\qquad m=N/p^{v_p(N)} .
\]
\end{proposition}

\begin{proof}
By \cref{prop:K5-reduction}, $A_N\equiv R(N)+\frac12(K_5(N)-Q(N))\pmod2$ for every
positive odd $N$ prime to $5$. Modulo $2$, $R(N)$ counts the representations
$N=5x^2+2y^2$ and $N=x^2+10y^2$ with $x$ odd and $x,y>0$. Since $\Leg{-10}p=-1$, a
representation with $p\mid N$ forces $p\mid x$ and $p\mid y$, and division by $p$ gives
$R(N)\equiv R(N/p^2)\pmod2$.

For odd $d$ prime to $5$, $q_0(d)\bmod4$ depends only on $d\bmod40$, and evaluating
the $16$ unit classes gives $q_0(d)+q_0(29d)\equiv2\pmod4$. Since
$p^2\equiv1\pmod{40}$, the two top $p$-power layers of divisors give
$Q(N)-Q(N/p^2)\equiv2\tau(m)\pmod4$.

Subtract the two formulas for $A$ and apply \cref{thm:K5-recurrence}. As $5\nmid m$,
$d_5(m)=\tau(m)$, these terms cancel, and $\nu_5(m)$ remains.
\end{proof}

\begin{theorem}\label{thm:gamma-scaling}
For every positive odd integer $r$ prime to $5$,
\[
  \gamma(29r)+\gamma(r)\equiv\begin{cases}0\pmod2,&r\equiv1,4\pmod5,\\1\pmod2,&r\equiv2,3\pmod5.\end{cases}
\]
\end{theorem}

\begin{proof}
Write $r=p^jm$ with $p\nmid m$. M\"obius inversion in \cref{prop:divisor-transform}
gives the exact identities
\begin{align*}
  \gamma(pm)+\gamma(m)&=\sum_{d\mid m}\mu(d)A_{pm/d},\\
  \gamma(p^{j+1}m)+\gamma(p^jm)&=\sum_{d\mid m}\mu(d)\bigl(A_{p^{j+1}m/d}-A_{p^{j-1}m/d}\bigr)\qquad(j\ge1);
\end{align*}
for the first, split the divisors of $pm$ into $d$ and $pd$ with $d\mid m$, and
for the second add the formulas for $\gamma(p^{j+1}m)$ and $\gamma(p^jm)$, which
telescope.
By \cref{prop:A-recurrence} either right side is $\sum_{d\mid m}\mu(d)\nu_5(m/d)$
modulo $2$. By definition $\nu_5=\one*\one_{\{2,3\bmod5\}}$, so M\"obius inversion
identifies this sum with $\one_{m\equiv2,3\ (5)}$. Multiplication by $29$ preserves the
set of nonresidues modulo $5$, so this equals $\one_{r\equiv2,3\ (5)}$.
\end{proof}

\begin{proof}[Proof of \cref{thm:main-scaling}]
\Cref{cor:scaling-nonsquare} proves each congruence when the associated number
$40n+(2\ell+1)^2$ is not a square. Suppose $40n+1=r^2$ with $r>0$. Then
$r\equiv1,9\pmod{10}$, so \cref{thm:gamma-scaling} gives
$C((29r)^2)+C(r^2)\equiv0\pmod2$. The substitution in the proof of
\cref{cor:scaling-nonsquare}, applied to \eqref{eq:hecke-rogers}, also gives
$h_0(841n+21)=-h_0(n)$ at square discriminants. \Cref{prop:reduction} now proves
\eqref{eq:scaling-b}.

If $40n+9=r^2$, then $r\equiv3,7\pmod{10}$, and \cref{thm:gamma-scaling} gives
$C((29r)^2)+C(r^2)\equiv1\pmod2$; again $h_1(841n+189)=-h_1(n)$. Write
$u(n)=\coef nU$. The exponents of $U$ are the $n$ with $40n+9=(10a+7)^2$, $a\ge0$.
Multiplication by $29$ exchanges the classes $3$ and $7$ modulo $10$, so exactly one of
$r$, $29r$ is $\equiv7\pmod{10}$ and $u(841n+189)+u(n)=1$. So the two correction terms
$2C_1+2U$ in $B_1\equiv H_1+2C_1+2U\pmod4$ together contribute $2\cdot1+2\cdot1\equiv0$
to the sum, and \eqref{eq:scaling-b1} follows.
\end{proof}

\section{Running the verifications}\label{sec:verification}

The data and code for the two computer-assisted steps, \cref{lem:cancellation}
and the Sturm-bound computation in the proof of \cref{thm:K5-recurrence}, are
described in \cref{app:certificate,app:verifier,app:sturm} and are available in
the repository
\begin{center}
\url{https://github.com/ramanujamath/non-rascoe-partitions-mod4}
\end{center}
in the directory \texttt{verification/appendix}. Download the three programs into one directory:
\begin{itemize}
\item \texttt{decode\_certificate.py},
\item \texttt{verify\_original\_certificate.py},
\item \texttt{prove\_square\_scaling\_29.py},
\end{itemize}
and run them, with Python $3.10$ or later and without the optimization flag
\texttt{-O}:
\begin{verbatim}
python3 decode_certificate.py
python3 verify_original_certificate.py
python3 prove_square_scaling_29.py
\end{verbatim}
Only the Python standard library is needed. The first program reconstructs the
certificate and checks its SHA-256 digest; the second reconstructs the
signatures of \cref{lem:cancellation} and checks every required identity (about
seven seconds); the third computes $\mathfrak F$ modulo $2$ through the Sturm
bound (about one second). Neither verification uses the search procedure that
produced the certificate, or any external data.


\appendix
\section{The cancellation certificate}\label{app:certificate}

The following program determines all nonzero values of the four functions of
\cref{sec:certificate}. A record has the form $[t,[b,a],[x,y,z],g]$ with
$t\in\{\texttt J,\texttt S,\texttt T,\texttt V\}$ naming $f_{\Jk},f_{\Sw},f_2,f_5$,
$(b,a)\in\mathcal B$, $0\le x,y,z<20$ and $g\in\{0,1\}$.

For a compact encoding, let $j\in\{0,1,2,3\}$ be the position of $t$ in
$(\texttt J,\texttt S,\texttt T,\texttt V)$ and $k\in\{0,\dots,6\}$ the position of $a$ in
$(0,1,2,3,4,5,7)$, counted from $0$, and assign to the record the integer
\[
  2\bigl(20\bigl(20\bigl(20\,(7(5j+b)+k)+x\bigr)+y\bigr)+z\bigr)+g .
\]
Sort these integers, take consecutive differences starting from $0$, and write
each difference in base $128$, least significant digit first, adding $128$ to
every digit except the last. The string \texttt{DATA} is the Base64 encoding of
the zlib-compressed byte sequence. The program inverts these steps and checks
the SHA-256 digest of the result.

{\footnotesize\VerbatimInput{code/decode_certificate.py}}

\section{The cancellation verifier}\label{app:verifier}

This program enumerates the possible signatures by exact Fourier--Motzkin
elimination and checks \eqref{eq:cancellation} on every combination of residues,
signature and bit (\cref{lem:cancellation}). It also checks the properties of
Jackson's map stated in \cref{sec:involution}. It reads the file produced by
\cref{app:certificate}.

{\footnotesize\VerbatimInput{code/verify_original_certificate.py}}

\section{The Sturm-bound verifier}\label{app:sturm}

This program computes $\mathfrak F$ modulo $2$ (\cref{sec:recurrence}) and checks
every coefficient of $U_{29}\mathfrak F$ required by Sturm's theorem. It also checks
the order of the character image used in the index computation. All polynomial
arithmetic over $\F_2$ is done exactly with integer bit operations. (The
comments in the program refer to the form $\mathfrak F$ as \texttt{F}.)

{\footnotesize\VerbatimInput{code/prove_square_scaling_29.py}}


\section*{Statements and declarations}
\noindent\textbf{Funding.} No external funding was received for this work.

\noindent\textbf{Competing interests.} The author declares no conflict of interest.

\noindent\textbf{Data availability.} All certificate data and verification code
needed for the computer-assisted parts of the proofs are included in the
appendices and in the accompanying repository.

\noindent\textbf{Author contribution.} Sumit Kumar Jha is the sole author and is
responsible for the content of the manuscript.

\noindent\textbf{Use of generative AI.} Generative AI tools were used to assist
with mathematical exploration, computational verification and manuscript
preparation. The author is responsible for the final content.

\begin{thebibliography}{99}

\bibitem{AAMW}
A. Alaca, \c{S}. Alaca, E. McAfee and K.~S. Williams,
\emph{Lambert series and Liouville's identities},
Dissertationes Math. \textbf{445} (2007), 1--72.

\bibitem{Andrews}
G.~E. Andrews, \emph{The fifth and seventh order mock theta functions},
Trans. Amer. Math. Soc. \textbf{293} (1986), 113--134.

\bibitem{Buzzard}
K. Buzzard, \emph{Notes on modular forms of half-integral weight},
online notes, 2 October 2013.

\bibitem{CGS}
S.-P. Cui, N.~S.~S. Gu and C.-Y. Su,
\emph{Generalizations of mock theta functions and radial limits},
Proc. Amer. Math. Soc. \textbf{151} (2023), 3317--3329.

\bibitem{Cuthbertson}
P. Cuthbertson, \emph{A requested analytic proof of an identity of Dixit, Kumar, and Srivastava},
arXiv:2608.30180v1 (2026).

\bibitem{DKS}
A. Dixit, G. Kumar and A. Srivastava,
\emph{Non-Rascoe partitions and a rank parity function associated to the Rogers--Ramanujan partitions},
arXiv:2508.04359v1 (2025).

\bibitem{Jackson}
T. Jackson, \emph{A short proof that every prime $p\equiv3\pmod8$ is of the form $x^2+2y^2$},
Amer. Math. Monthly \textbf{107} (2000), 447.

\bibitem{Sills}
A.~V. Sills, \emph{$q$-Difference equations and identities of the Rogers--Ramanujan--Bailey type},
J. Difference Equ. Appl. \textbf{10} (2004), 1069--1084.

\bibitem{Stein}
W.~A. Stein, \emph{Modular Forms, a Computational Approach},
Graduate Studies in Mathematics \textbf{79}, Amer. Math. Soc., Providence, 2007.

\end{thebibliography}


\end{document}
